% ============================================
% 第15章：请把手弄脏——关于极致野望的七个字
% 组合桥 · 终章
% ============================================

\chapter{请把手弄脏}

\vspace{0.5cm}

\begin{center}
{\small\color{muted} 这座桥，来自2025—2026年的课堂、项目、录音与反复修改。}
{\small\color{muted} 所有的公式、所有的故事、所有的失败和骄傲，都指向这五个字。}
\end{center}

\vspace{1cm}

\section{野望}

野望不是``我要考第一''。不是``我要拿奖''。不是``我要比别人强''。

野望是：\textbf{``我要让这座桥站住。''}

2025年10月21日，我在结构设计大赛的赛场上。一个学生做了一个100克的船，载了19公斤。我在课堂上讲这个故事的时候，说``也不代表啥''。但语气里，有一种藏不住的骄傲。不是骄傲排名。是骄傲\textbf{那个东西站住了}。

2026年7月1日，学生在结题汇报里谈到AI画错的弯矩图，也谈到版本覆盖与反复修改。他最后说：``你就真的得是自己会才行，要不然你也不知道AI做得对不对。''这句话没有口号那么整齐，却更接近项目的真相：\textbf{工具可以推进工作，判断不能凭空生成。}

这就是野望。不是超过别人。是\textbf{做成一件自己本来做不成的事}。

\begin{insight}
\label{insight:chap15}

\textbf{野望不是竞争。野望是创造。}

不是``我要比别人强''。是``我要亲手做出一个东西，它以前不存在，现在存在了，而且它站住了''。

这个世界上，最稀缺的不是聪明人。聪明人很多。最稀缺的是\textbf{愿意把手弄脏的人}——愿意花三个小时搭一座桥，塌了五次，第六次站住了，然后抬起头，眼睛里有光，说：``你看，它站住了。''
\end{insight}

\section{技艺的幸福感}

2026年6月，SpanCraft设计复盘。我写下了那八个维度。但八个维度背后，有一个更深的东西——我后来才意识到，那才是SpanCraft真正在做的事。

\textbf{SpanCraft在教学生感受``技艺的幸福感''。}

什么是技艺的幸福感？是你亲手做了一样东西，你知道它每一个细节，你为它花了时间，你为它失败过，你为它骄傲过。你不需要别人告诉你它好不好。你自己知道。你看着它，心里有一种\textbf{安静}的满足。

这种满足，AI给不了你。AI可以生成一座完美的桥，比你搭的漂亮一百倍。但那是AI的桥。不是你的桥。你看着AI的桥，不会心跳加速。你看着\textbf{你自己搭的桥}——那座歪歪扭扭、梁太高、索太粗、但它站住了的桥——你才会心跳加速。

\textbf{因为那座桥上有你的时间。你的错误。你的决定。你的骄傲。}

\section{五个字}

这本书叫《请把手弄脏》。五个字。

这不是一个比喻。我真的是在说：请伸出你的手。去碰混凝土。去弯钢筋。去调参数。去写代码。去搭一座桥。然后看它塌掉。再搭一座。再看它塌掉。再搭一座——直到它站住。

\textbf{不是AI做不了。是AI做了，就不是你的了。}

2026年5月22日，我在游戏化教学分享中说了那句话：``学生在学习过程中从来没有失败的权利，他也就没有成功的那个喜悦或者荣耀。''这句话，被反复引用，被写进了教学理念分析，被印在了这本书的第2章。

但我想在这里再说一遍。因为它是全书最重要的一句话。

\textbf{失败的权利，是学习的权利。}如果你从来没有失败过——如果你从来没有搭过一座会塌的桥，从来没有写过一段会报错的代码，从来没有问过一个``蠢''问题——那你从来没有真正学过任何东西。你只是被灌输过。被灌输的东西，AI可以做得更好。但你自己学到的东西——那些在你失败过、调整过、最终理解了的东西——AI永远拿不走。

\section{桥的本事}

桥的本事，是把原本分开的两岸，连成彼此的两岸。

海德格尔说这句话的时候，他说的不是工程。他说的是\textbf{存在}。一座桥不是简单地连接了A点和B点。一座桥让A点成了``此岸''，让B点成了``彼岸''。在桥出现之前，两岸只是两块土地。在桥出现之后，它们成了彼此的对方。

\textbf{一堂好课，做的也是同一件事。}

在好的课堂出现之前，学生和学生是两个互不相干的个体。在好的课堂出现之后，他们成了``我们''——一群一起搭过桥的人，一起失败过、一起骄傲过、一起在深夜调过参数的人。

\textbf{一个好游戏，做的也是同一件事。}

在好游戏出现之前，``知识''和``我''是两个互不相干的东西。在好游戏出现之后，知识成了我手指上的感觉——那个$qL^2/8$，不再是黑板上的一行公式，而是我亲手让它塌过三次之后，终于理解了的东西。

\textbf{一本好书，做的也是同一件事。}

在好书出现之前，``AI时代的教学''是一个让人焦虑的话题。在好书出现之后——我希望——它变成了一个让人兴奋的邀请。不是``你应该学''。不是``你应该改''。是\textbf{``你想不想来试试？''}

\section{在一个略显荒凉的时代里}

这本书写于2025—2026年。

这不是一个对土木工程特别友好的时代。分数线在掉。新生群里有人劝退。AI可以生成一切``看起来像桥''的东西。

但这也是一个对土木工程特别需要的时代。因为\textbf{当AI能生成一切的时候，真正稀缺的，就不再是生成，而是判断。不再是速度，而是责任。不再是看起来像，而是真的站得住。}

土木工程教会我的，恰好就是这些。

\textbf{你不能骗一座桥。}你可以骗一份PPT。你可以骗一份报告。你可以骗一个评委。你可以骗你自己。但你不能骗一座桥。桥不会听你解释。它只会站住，或者塌掉。

在一个越来越多人相信``差不多就行''的世界里，这种``不能骗''的品格，是一种稀缺的、值得被保护的东西。

\textbf{创造美，并且严肃地为这份美负责。}这大概是工程教育里最难教、也最值得教的东西。而要教会它，第一步，是先让人重新愿意\textbf{把手弄脏}。

\section{带着骄傲，和一点悲伤}

2026年6月，SpanCraft设计复盘。结尾我写了这样一段话：

\begin{shiyu-inline}
``桥的本事，是把分开的两岸连起来。也许一堂好课、一个好游戏，做的也是同一件事。在一个略显荒凉的时代里，把'骄傲'，和那些'还愿意相信的人'，重新连接起来。带着骄傲，和一点悲伤，我把这座桥，建给每一个愿意从它背上走过的人。''

\hfill ——SpanCraft设计复盘，2026年6月
\end{shiyu-inline}

\textbf{带着骄傲，和一点悲伤。}

骄傲，是因为我们做到了。我们证明了：好玩和有料，不矛盾。AI和人的判断，不矛盾。效率和意义，不矛盾。我们证明了：在AI能做一切的时代，教育仍然可以是一件让人心跳加速的事。

悲伤，是因为不是所有人都看到了。还有无数学生在问``学这个有什么用''。还有无数老师在担心``AI会不会取代我''。还有无数课堂，仍然把大量时间用在复述那些工具已经可以迅速生成的内容上。

但这本书，就是一座桥。

让一个从没想过土木的人，因为它好看好玩，愿意多停留三十秒。让一个正在犹豫要不要转专业的学生，忽然想起来——结构，原来是有美的。让一个被AI吓到的教师，重新发现——\textbf{工具能力在变，但人仍要决定什么值得教、如何核验，以及由谁负责}。

\begin{center}
{\LARGE\bfseries 请把手弄脏。}
\end{center}

\vspace{2cm}

\begin{flushright}
ColdCat（大连交通大学）\\
2026年7月，于大连
\end{flushright}
